\chapter{Background and Related Work}\label{ch:background}

\section{Learning manipulation from demonstrations}

In imitation learning a policy $\pi$ maps an observation $o_t$ (camera images and the robot's joint positions) to an
action $a_t$ (target joint positions) and is trained on demonstrations recorded by a human. With low-cost arms such as
the SO-101~\cite{so101}, demonstrations are recorded by teleoperation: the operator moves a \emph{leader} arm by hand
and an identical \emph{follower} arm copies it, while cameras and joint encoders are recorded at a fixed rate. LeRobot
stores such recordings as datasets of episodes and trains policies on them~\cite{lerobot}.

\section{Action chunking}

Predicting one action per step makes policies sensitive to compounding errors and to the timing of inference.
\emph{Action chunking} predicts a sequence of $H$ future actions $A_t=(a_t,\dots,a_{t+H-1})$ from one observation.
ACT~\cite{zhao2023act} introduced chunking with a transformer conditional variational autoencoder. Diffusion
Policy~\cite{chi2023diffusion} generates the chunk by iterative denoising. Most current VLAs generate chunks with a
\emph{flow-matching} action head conditioned on the backbone's features: \smolvla~\cite{shukor2025smolvla} pairs a
small vision-language model with an action expert, GR00T~N1.5 and~N1.7~\cite{nvidia2025gr00t} use a diffusion
transformer over a 3-billion-parameter model, and \pizerofive~\cite{pi2025pi05} a larger flow-matching expert.
MolmoAct~2~\cite{molmoact2} adds reasoning about actions in space. The chunk length differs widely: 16 actions for
GR00T, 50 for \smolvla{} and \pizerofive, 100 for ACT, which matters for remote inference (Section~\ref{sec:shortchunk}).

\section{Asynchronous and remote inference}\label{sec:bg-remote}

If the robot waits for each chunk before moving, it pauses for every inference. \emph{Asynchronous} inference requests
the next chunk while the current one is still being executed, so execution and inference overlap; \smolvla{}'s authors
proposed this in a client-server form~\cite{shukor2025smolvla}. LeRobot~0.7 ships a remote-inference stack built on this
idea: a policy server on the GPU machine and a client on the robot computer, connected through the Zenoh protocol.
The client keeps a queue of actions, has at most one request in flight, discards replies that arrive after their
deadline, and requests the next chunk when the queued actions fall to about one round trip. If no action is available
for a configured \emph{starvation grace} (one second by default), it ends the run with a fault.

\section{Real-time chunking}

Consecutive chunks are predicted from different observations, so the actions at the boundary between them can
disagree and the arm jumps when execution switches chunks. Real-Time Chunking (RTC)~\cite{black2025rtc} treats the
next chunk as an in-painting problem: while generating it with a flow-matching policy, the actions that will still be
executed from the current chunk are fixed as a prefix, so the new chunk continues them smoothly. RTC needs a
flow-matching policy and an estimate of how many actions will be executed before the new chunk arrives.

\section{Distributed Real-Time Chunking}\label{sec:bg-drtc}

\drtc~\cite{vial2026drtc} combines asynchronous inference with RTC and makes the client resilient to an unreliable
network:
\begin{itemize}
\item \textbf{Trigger.} A new chunk is requested when the queue holds at most $H-s_{\min}$ actions, where $H$ is the
      chunk length and $s_{\min}$ a fixed number of steps.
\item \textbf{Latency estimation.} The round-trip time is estimated with the Jacobson--Karels
      estimator~\cite{jacobson1988congestion,rfc6298} that TCP uses for its retransmission timer: a smoothed mean
      plus a multiple of a smoothed mean deviation.
\item \textbf{Cooldown.} A request that is not answered within a cooldown derived from the latency estimate is given
      up, and a fresh observation is sent instead; the freshest observation always wins.
\item \textbf{Message semantics.} Messages between threads and the merging of chunks into the action schedule are
      modelled as last-write-wins registers from the literature on conflict-free replicated data
      types~\cite{shapiro2011crdt}, so reordered and duplicated messages are absorbed.
\end{itemize}
Its author evaluated it with injected faults on an SO-101 arm. The present work ports its scheduling into LeRobot~0.7's
runtime and evaluates it against the shipped client on task success.

\section{Smoothing robot commands}

Robot commands that change abruptly excite the arm's mechanics and look unnatural. A standard remedy is a causal
low-pass filter on the commanded positions; the Butterworth filter~\cite{butterworth1930} has a maximally flat pass
band. A second remedy, specific to chunked control, is to vary the playback rate: when the next chunk will arrive late,
executing the queued actions more slowly avoids a full stop.

\section{Summary}

Remote inference makes large policies usable on low-cost robots, but the shipped client gives up quickly under network
faults. \drtc{} proposes a scheduling that tolerates such faults. What is missing is an independent, controlled
comparison on a physical robot, against the shipped client, with task success as the outcome and with the motion
quality measured as well. This report provides one.
